(1) 外侧自由传播波数为 $k=\sqrt{2mE}/\hbar$，势垒内衰减常数为 $\kappa=\sqrt{2m(V_0-E)}/\hbar$。在 $x<0$ 写入射与反射波，在 $0<x<a$ 写衰减和增长两指数的线性组合，在 $x>a$ 只保留右行透射波。由于势垒是有限跳跃势，$\psi$ 和 $\psi'$ 在 $x=0,a$ 两界面都连续，不能只保留单一衰减指数来满足两个边界。
(2) 四个匹配条件消去振幅后得到精确的 $T(E)=[1+V_0^2\sinh^2(\kappa a)/(4E(V_0-E))]^{-1}$。在 $\kappa a\gg1$ 下 $\sinh^2(\kappa a)\simeq e^{2\kappa a}/4$，故 $T\simeq16E(V_0-E)e^{-2\kappa a}/V_0^2$。其中指数项是隧穿穿越势垒的主抑制，前因子来自边界匹配，不能简单宣称精确 $T=e^{-2\kappa a}$。
(3) 若 $E>V_0$，势垒内也是传播波，其波数 $q=\sqrt{2m(E-V_0)}/\hbar$。解析延拓 $\kappa\to iq$ 后得到 $T(E)=[1+V_0^2\sin^2(qa)/(4E(E-V_0))]^{-1}$。当 $qa=n\pi$、$n=1,2,\ldots$ 时内部往返相位恰好匹配，$\sin(qa)=0$，透射率为一；这是有限势垒内多次反射相消导致的共振透射，而非经典粒子完全不受势垒影响。
(4) 在阈值 $E\to V_0$ 时，利用 $\sinh(\kappa a)\sim\kappa a$ 或 $\sin(qa)\sim qa$，两侧表达式都趋于 $T(V_0)=[1+mV_0a^2/(2\hbar^2)]^{-1}$。此处 $\kappa\to0$，条件 $\kappa a\gg1$ 被破坏，不能把厚势垒的指数渐近式带入。势垒顶端的真实波函数是线性函数的极限，而不是某个仍快速衰减的指数。
